\begin{abstract}
To answer a question, one must first determine what it asks. Questions about videos often admit several readings, since they may refer to multiple entities or moments. While humans respond to such referential ambiguity by seeking clarification or addressing each alternative, most video-language models commit to a single reading, a behavior we term \textit{hallucinated commitment}. We examine this failure by introducing \textbf{\bench{}}, a benchmark that pairs each question with exhaustive, human-validated readings and per-reading answers. Unlike in images, a video question's readings are distributed over time. Most models commit silently even when invited to ask. The models that do ask often answer the named reading incorrectly. Fine-tuning induces \textit{format mimicry}: models adopt the style of ambiguity-aware answers but fill them with incorrect readings. Once the validated readings are supplied, the open-weight models answer nearly perfectly, so for them the bottleneck is finding the readings, not answering them. Our results underscore the urgency of equipping models to determine what a question asks before answering it.
\end{abstract}
